% Lösungen Grenzwerte und Verhalten im Unendlichen (zusammenbau v0.4, Heft)
\einheitenkopf{Grenzwerte und Verhalten im Unendlichen}

\erg{55}{a) Grad ungerade, Koeffizient negativ \quad b) für $x \to +\infty$: $f(x) \to +\infty$; für $x \to -\infty$: $f(x) \to +\infty$}
\erg{56}{a) für $x \to +\infty$: $f(x) \to -\infty$; für $x \to -\infty$: $f(x) \to +\infty$ \quad b) für $x \to +\infty$: $f'(x) \to +\infty$; für $x \to -\infty$: $f'(x) \to +\infty$ \quad c) für $x \to +\infty$: $f'(x) \to -\infty$; für $x \to -\infty$: $f'(x) \to +\infty$ \quad d) für $x \to +\infty$: $f'(x) \to -\infty$; für $x \to -\infty$: $f'(x) \to +\infty$}
\erg{57}{a) für $x \to -\infty$ \quad b) für $x \to +\infty$: $f(x) \to +\infty$; für $x \to -\infty$: $f(x) \to 0$}
\erg{58}{a) für $x \to +\infty$: $f(x) \to -\infty$; für $x \to -\infty$: $f(x) \to 0$, Annäherung von oben \quad b) für $x \to +\infty$: $f(x) \to +\infty$; für $x \to -\infty$: $f(x) \to 0$, Annäherung von unten \quad c) für $x \to +\infty$: $f(x) \to -\infty$; für $x \to -\infty$: $f(x) \to 0$, Annäherung von oben \quad d) $f(x) \to 0$: der e-Faktor fällt schneller, als $x + 4$ wächst; der Graph nähert sich von oben der $x$-Achse. \quad e) $f(x) \to 0$: die e-Funktion ist stärker als der wachsende Faktor; der Graph nähert sich von oben der $x$-Achse. \quad f) $f(x) \to 0$: der e-Faktor fällt schneller, als der lineare Faktor wächst; der Graph nähert sich von oben der $x$-Achse. \quad g) für $x \to +\infty$: $f(x) \to +\infty$; für $x \to -\infty$: $f(x) \to 0$ \quad h) für $x \to +\infty$: $f(x) \to 0$; für $x \to -\infty$: $f(x) \to -\infty$ \quad i) $f(x) \to 0$: $e^{x}$ geht gegen null und gewinnt gegen das Polynom; $x \cdot (4 - x)$ ist dort negativ, der Graph nähert sich von unten der $x$-Achse. \quad j) $w(t) \to 40$: das Produkt geht gegen null, der Summand bleibt; der Pegel sinkt auf Dauer wieder auf 40 cm. \quad k) $b(t) \to 90$: das Produkt geht gegen null, der Summand bleibt; der Blutzucker kehrt auf Dauer zum Nüchternwert 90 mg/dl zurück. \quad l) $T(t) \to 18$: das Produkt geht gegen null, der Summand bleibt; das Rohr kühlt auf Dauer auf 18 °C ab.}
\erg{59}{a) Nullstelle $x = 5$ (der e-Faktor ist nie null); für $x \to +\infty$: $f(x) \to +\infty$; für $x \to -\infty$: $f(x) \to 0$, von unten \quad b) Nullstelle $x = 3$ (der e-Faktor ist nie null); für $x \to +\infty$: $f(x) \to -\infty$; für $x \to -\infty$: $f(x) \to 0$, von oben \quad c) Nullstelle $x = -7$ (der e-Faktor ist nie null); für $x \to +\infty$: $f(x) \to 0$; für $x \to -\infty$: $f(x) \to -\infty$}
